\firstsection{Introduction}
\maketitle
\raggedbottom
\setlength{\emergencystretch}{1.5em}
Time-varying scalar fields record how quantities such as sea-level pressure, aerosol concentration or vorticity evolve over a spatial domain. Analysts study them through features, such as low-pressure systems, pollution episodes or shed vortices, and ask three questions about each: how it relates to the others, where it is, and how it moves. Animation keeps every time step in its spatial context but makes distant times hard to compare~\cite{Tversky:2002:Animation,Robertson:2008:Animation}. Static overviews therefore assign one display axis to time and compress space into the other, so that a whole sequence fits into one image~\cite{Bach:2017:STC}. What that second axis means decides which of the three questions the overview can answer. The Hovm\"oller diagram of atmospheric science keeps a physical coordinate such as longitude and shows an eastward-moving system as a slanted streak~\cite{Hovmoller:1949}, but it averages the field over the remaining coordinates, so neighboring features blend and their hierarchy disappears.

Merge tree maps answer the first question. A merge tree records how sublevel sets connect as the threshold rises: its leaves are the extrema, and its internal nodes are the levels at which regions merge~\cite{Heine:2016:Survey,Edelsbrunner:2002:TPS}. Temporal merge tree maps (TMTM) draw the tree of each time step as a one-dimensional profile whose minima are the extrema and whose barriers are the merge levels, and they stack the profiles into one image~\cite{Koepp:2023:TMTM}. Spatiotemporal merge tree maps (ST-MTM) additionally arrange the features by their pairwise distances and temporal coherence~\cite{STMTM:2026}. In both, the vertical axis follows the tree, temporal alignment or relative distance; neither fixes an origin or a direction in space. A storm that crosses the North Atlantic within a few days can therefore stay at the same height of the map, and the reader sees how it connects to neighboring lows but not where it is or which way it moves.

We give this axis a spatial meaning. Fixing one direction $a^*$ in the domain assigns every extremum a reference position, the projection of its location onto $a^*$, and a map whose axis encodes location should draw each extremum at its reference (Fig.~\ref{fig:storms_overview}). The obstacle is the tree itself. A profile shows the tree exactly only if every subtree occupies one contiguous block of leaves, but the references of two sibling subtrees may interleave (Fig.~\ref{fig:related-work-mapping}d). No exact profile can then place every feature at its reference, whatever algorithm draws it. This conflict raises three questions: how much displacement the tree forces, where in the tree this displacement arises, and how much merge-level fidelity must be given up to remove it. Without answers, the positional error of a map cannot be attributed to the data or to its layout, and an analyst cannot judge whether a spatial axis is worth its cost.

We answer these questions with bounds and turn the answers into a construction. Three lower bounds, for exact topology, for features of finite width and for a bounded merge error, depend only on the trees, the references and the widths, so the positional price of the axis is known before any layout is computed and applies to every map that keeps the tree. AnchorMTM then computes the map directly from the field sequence and one merge budget $\varepsilon$. It searches leaf orders that the budget admits, using a closed-form test for the smallest merge error an order can show, places the extrema as close to their references as the chosen order allows, and reconstructs the profiles. With $\varepsilon=0$, the map shows every tree exactly; with $\varepsilon>0$, it releases the conflicts that a small change of merge levels can resolve and records the merge error each frame actually spends.

Our contributions are as follows.
\begin{itemize}\setlength{\itemsep}{1pt}\setlength{\parsep}{0pt}
\item We formulate spatially anchored merge tree maps and derive three layout-independent lower bounds on the displacement that keeping the tree forces, for exact topology, finite widths and bounded merge error. The bounds locate the conflict in the tree and measure any merge tree map, including existing ones.
\item We present AnchorMTM, which builds the leaf order and extremum positions of every frame from the projected locations, so that the height of a band reads as location and its drift as motion along $a^*$, while the profile keeps the nesting and merge levels of the tree. A merge budget $\varepsilon$, enforced through optimal merge barriers and checked on every rendered frame, controls how much merge-level fidelity may be traded for position.
\item On five planar and two volumetric sequences, we show that the bounds are tight on controlled and real trees, that most of the order error of existing maps is not forced by the tree, and how a budget reshapes the displacement of individual frames.
\end{itemize}
